\input{../tex-inputs/latex-leseansicht-vorspann}

\section[Arthur Schnitzler an Gustav Schwarzkopf, 7. 8. 1921]{L04463 Arthur Schnitzler an Gustav Schwarzkopf, 7. 8. 1921}
\nopagebreak\mylabel{L04463v}
\rehead{ }\normalsize\beginnumbering\briefempfaengerindex{Schwarzkopf, Gustav@\textsc{Schwarzkopf, Gustav}!zzzSchnitzler, Arthur@\emph{von Arthur Schnitzler}!1921-08-071@{7. 8. 1921}|(be}
\toendnotes[C]{\smallbreak\pagebreak[2]}
\correspDesc{Versand  durch Arthur Schnitzler am 7. 8. 1921 in Altaussee
\newline{}Übermittlung  am 8. 8. 1921 in Altaussee
\newline{}Erhalt  durch Gustav Schwarzkopf im Zeitraum [8. 8. 1921
                  – 12. 8. 1921?] in Wien}\toendnotes[C]{\smallbreak}
\Standort{DLA, A:Schnitzler, HS.1985.1.1897.}
\physDesc{Bildpostkarte, 496 Zeichen
\newline{}Handschrift: Bleistift, lateinische Kurrent
\newline{}Versand: Stempel: »\nobreak{}\oindex{Altaussee@\textbf{Altaussee}|pwk}Alt Aussee, 8. VIII. 21\nobreak{}«.  }\toendnotes[C]{\smallbreak}\pstart{}{\pb}Hrn Gustav Schwarzkopf\pend{}\pstart{}Wien\oindex{Wien@\textbf{Wien}|pw}\pend{}\pstart{}I. Tiefer Graben 17\oindex{Wien@\textbf{Wien}!I., Innere Stadt@\textbf{I., Innere Stadt}!Tiefer Graben 17@\textbf{Tiefer Graben 17}|pw}\pend{}{\bigskip}
\pstart
           \noindent{}\centering{}{\pb}\textcolor{gray}{\textbf{Salzkammergut. Alt-Aussee\oindex{Altaussee@\textbf{Altaussee}|pw}.}}\pend
           
\pstart
           {[}hs. Schnitzler:{]} \label{T_L04463-1v}\edtext{Mein Fenster}{\lemma{\textnormal{\emph{Mein Fenster}}}\Cendnote{\textnormal{Schnitzler markiert sein Zimmer im Hotel Seewirt\oindex{Hotel am See@\textbf{Hotel am See}|pwk} mit einer Bleistiftlinie.}}}\label{T_L04463-1}\pend
           \vspace{1em}
\pstart
           \raggedleft{}{\pb}7. 8. 1921\pend
           \vspace{0.5em}
\pstart
           lieber Gustav,{ }\label{K_L04463-1v}\edtext{seit Dienstg 27}{\lemma{\textnormal{\emph{seit Dienstg 27}}}\Cendnote{\textnormal{Schnitzler war am Dienstag, den 26. 7. 1921
                  angereist, vgl. A. S.: \emph{Tagebuch}, 26. 7. 1921.}}}\label{K_L04463-1} bin ich hier – ein Regentag bisher; sonst das wunderbarste Wetter; –
               ein heimisches Bade u Strandleben hat sich entfaltet. Die Kinder\pwindex{Cappellini, Lili 13.\,9.\,1909 Wien – 26.\,7.\,1928 Venedig@\textsc{Cappellini, Lili} (13.\,9.\,1909 Wien – 26.\,7.\,1928 Venedig)|pwv}\pwindex{Schnitzler, Heinrich 9.\,8.\,1902 Hinterbrühl – 12.\,7.\,1982 Wien@\textsc{Schnitzler, Heinrich} (9.\,8.\,1902 Hinterbrühl – 12.\,7.\,1982 Wien)|pwv} sind \label{K_L04463-2v}\edtext{seit
               31. 7.}{\lemma{\textnormal{\emph{seit
               31. 7.}}}\Cendnote{\textnormal{Auch hier vertat sich Schnitzler im Datum: Die Kinder\pwindex{Cappellini, Lili 13.\,9.\,1909 Wien – 26.\,7.\,1928 Venedig@\textsc{Cappellini, Lili} (13.\,9.\,1909 Wien – 26.\,7.\,1928 Venedig)|pwkv}\pwindex{Schnitzler, Heinrich 9.\,8.\,1902 Hinterbrühl – 12.\,7.\,1982 Wien@\textsc{Schnitzler, Heinrich} (9.\,8.\,1902 Hinterbrühl – 12.\,7.\,1982 Wien)|pwkv} reisten am 30. 7. 1921 zu ihrer Mutter\pwindex{Schnitzler, Olga 17.\,1.\,1882 Wien – 13.\,1.\,1970 Lugano@\textsc{Schnitzler, Olga} (17.\,1.\,1882 Wien – 13.\,1.\,1970 Lugano)|pwkv}, vgl. A. S.: \emph{Tagebuch}, 30. 7. 1921.}}}\label{K_L04463-2} in Possenhofen\oindex{Possenhofen@\textbf{Possenhofen}|pw}; ich
               bleibe wohl bis 25. vielleicht länger hier, da{\geminationn}{ }Bayern\oindex{Bayern@\textbf{Bayern}|pw}. Sehe viel, aber meist angenehme Bekannte;
               Geselligkeit {\pb}ohne Zwang. – Mein Zimmer\oindex{Hotel am See@\textbf{Hotel am See}|pwv} sehr schön; Landschaft entzückt mich
               immer mehr.\pend
           \pstart Sein Sie, mit mit Ihrem Bruder\pwindex{Schwarzkopf, Max 12.\,6.\,1857 Wien – 14.\,4.\,1928 ebd.@\textsc{Schwarzkopf, Max} (12.\,6.\,1857 Wien – 14.\,4.\,1928 ebd.)|pwv} herzlichst gegrüßt. Ihr \spacefill\mbox{A.}\pend{}\selectlanguage{ngerman}\endnumbering\briefempfaengerindex{Schwarzkopf, Gustav@\textsc{Schwarzkopf, Gustav}!zzzSchnitzler, Arthur@\emph{von Arthur Schnitzler}!1921-08-071@{7. 8. 1921}|)be}\mylabel{L04463h}
\begin{anhang}
\end{anhang}\newcommand{\dateiname}{L04463}\newcommand{\titel}{Arthur Schnitzler an Gustav Schwarzkopf, 7. 8. 1921}\newcommand{\editorInnen}{Herausgegeben von Jahnke, SelmaMüller, Martin Anton}\input{../tex-inputs/latex-leseansicht-abspann}
